\section{Methodology}

\subsection{Study Design}

This study uses an integrative evidence synthesis design combining a PRISMA 2020-guided systematic review~\cite{page2021prisma} with a structured comparative multi-case analysis of publicly documented national eHealth implementations. The systematic review provides a broad, synthesized evidence base of adoption determinants from the peer-reviewed and grey literature. The multi-case analysis applies the U-T-I-O framework to real-world implementation contexts, enabling validation and comparative assessment of how determinants manifest across different national and organizational settings. No primary data involving human subjects were collected or analyzed. Prospective protocol registration was not undertaken: the review's outcomes are implementation and adoption determinants rather than patient health outcomes, placing it outside PROSPERO's registration scope, and the synthesis draws exclusively on published literature and public documentation. To mitigate the absence of registration, the search strategy, inclusion and exclusion criteria, and U-T-I-O coding framework were defined a priori, before search execution, and are fully documented in Supplementary File~S1.

\subsection{Systematic Review Protocol}

\subsubsection{Database Selection and Search Strategy}

The systematic literature search reported here covers three primary databases—PubMed/MEDLINE, IEEE Xplore, and Scopus—supplemented by Web of Science, CINAHL, Google Scholar, hand-searching of reference lists, and expert recommendations, for publications from 2003 through May 2026. The 2003 lower bound was selected to incorporate the DeLone and McLean IS Success Model~\cite{delone2003model} as a foundational theoretical anchor; the bulk of the empirical evidence base is drawn from 2015 onwards, covering the HITECH incentive era through current FHIR mandate implementation. The primary databases capture interdisciplinary literature spanning biomedical informatics, computer science, and health policy; the supplementary sources broaden coverage into nursing/allied health, cross-disciplinary outlets, and grey literature. Each search combined terms for the technology domain (``electronic health record,'' ``EHR,'' ``health information exchange,'' ``eHealth,'' ``interoperability'') with adoption-related terms (``adoption,'' ``implementation,'' ``barriers,'' ``facilitators,'' ``determinants,'' ``success factors''), with database-specific field tags, MeSH terms, and document-type filters applied. Full Boolean query strings for each database are provided in Supplementary File~S1.

\subsubsection{Inclusion and Exclusion Criteria}

Records are included if they: (1) appear in peer-reviewed journals or as institutional/government reports; (2) focus on eHealth, EHR, EMR, HIE, health information technology, or clinical interoperability; (3) address adoption, implementation, barriers, facilitators, or determinants of system uptake; and (4) employ empirical, systematic review, scoping review, or framework-based research designs. Records are excluded if they: (1) are opinion editorials, letters, or conference abstracts without full methodology; (2) address non-healthcare IT systems without healthcare application; or (3) are published in languages other than English. Studies involving the primary collection of identifiable human-subject data are also excluded.

\subsubsection{Screening and Data Extraction}

Retrieved records are deduplicated using reference management software. Title and abstract screening is applied to all retained records using the inclusion and exclusion criteria. Full-text review is conducted for records passing the initial screen. Data extraction is conducted using the U-T-I-O coding framework, applying thematic codes corresponding to each of the four adoption domains across all included sources. The resulting evidence base is documented in the evidence extraction matrix (see supplementary materials).

\subsection{Multi-Case Analysis}

\subsubsection{Case Selection}

Four national eHealth implementations were selected for comparative analysis using purposive sampling~\cite{yin2018casestudy,patton2015qualitative,adams2020strategies}. Cases were selected to maximize variation in national context, governance structure, and adoption maturity, while ensuring that sufficient publicly available documentation existed for systematic analysis. Selected cases include:

\begin{itemize}
    \item \textbf{U.S. ONC Health Information Exchange Programs} — representing the dominant federated, market-based model of HIE development, supported by HITECH Act incentives~\cite{onc2025reports}
    \item \textbf{NHS Shared Care Records (UK)} — representing a centrally coordinated national interoperability program within a single-payer health system~\cite{nhs2025sharedcare}
    \item \textbf{Estonia National eHealth System} — representing a comprehensive, legally mandated national digital health infrastructure built on the X-Road data exchange platform~\cite{estonia2026ehealth}
    \item \textbf{U.S. VA/DoD Federal EHR Modernization} — representing a large-scale federal EHR replacement program with significant procurement and implementation challenges~\cite{gao2025vaehr,gao2024dodehr}
\end{itemize}

\subsubsection{Case Analysis Protocol}

Each case was analyzed using the U-T-I-O framework as the organizing construct. Primary sources—including official government reports, national agency publications, legislative documents, and independent audit findings—were reviewed and coded against the four domains. Domain-level ratings (High / Medium / Low--Medium / Low) were assigned via a preponderance-of-evidence rule using a four-construct rubric per domain; complete construct definitions, decision rules, and per-case justification are provided in Supplementary File~S3. To check for single-reviewer bias, a supplementary AI-assisted second-rater pass re-rated all 16 case-domain cells blind to the original ratings (protocol, per-cell justifications, and adjudication log in Supplementary File~S4). Within-one-band agreement was 15/16 (93.8\%); exact agreement was 6/16 (37.5\%), with all disagreements running in the conservative direction. This led to the revision of two VA/DoD cells—User Readiness and Technical Interoperability—from Medium to Low--Medium, aligning them with the rubric's explicit decision-rule triggers. Agreement statistics, reported in full in the Limitations section and Supplementary File~S4, reflect the pre-revision ratings. Independent dual-rating with a credentialed human second rater remains an outstanding limitation of this study.

\subsection{Analytical Model}

The U-T-I-O Model operationalizes eHealth adoption as a function of four interdependent domains:

\begin{equation}
Adoption = f(U, T, I, O)
\end{equation}

where U denotes User Readiness, T denotes Technical Interoperability, I denotes Institutional Alignment, and O denotes Organizational Execution. The model posits that sustained adoption is most likely when all four domains achieve sufficient alignment, and that weakness in any single domain can constrain the adoption value realized even when other domains are strong.
